\chapter{Cryptographic Subspaces over Finite Fields}

\section{From the Rationals to $\mathbb{F}_2$}

While the initial exploration of the $M_{47}$ tensor deformation occurred over the rational numbers $\mathbb{Q}$, its true power manifests when generalized to arbitrary fields, specifically the finite field of two elements, $\mathbb{Z}/2\mathbb{Z}$ (or $\mathbb{F}_2$).

The coefficients defining the $M_{47}$ deformation matrix:
$$ \begin{bmatrix}
1 & \frac{13}{19} & \dots \\
\frac{3}{7} & 1 & \dots
\end{bmatrix} $$
were refactored to use formal field inverses (e.g., $3 \cdot 7^{-1}$). Provided the characteristic of the field does not divide the denominator, the tensor remains well-defined. Over $\mathbb{F}_2$, these topological structures warp into discrete Boolean manifolds.

\section{Applications to Module-LWE}

The Module Learning With Errors (Module-LWE) problem forms the bedrock of post-quantum cryptography. The problem requires identifying a hidden vector $s$ given samples of the form $(a, a \cdot s + e) \in (R_q)^2$, where $R_q = \mathbb{Z}_q[X]/(X^N+1)$.

By injecting the $M_{47}$ fractional topology into the coefficient rings, the AI framework demonstrated that the lattice basis reduction algorithms (like LLL and BKZ) encounter non-Euclidean "potholes"—subspaces where the volume heuristics break down due to the intersection metric $\chi^{5/2}$.

\section{The Future of Automated Mathematical Discovery}

The ability to port complex algebraic structures seamlessly from $\mathbb{R}$ to $\mathbb{F}_2$ within a fully verified Lean 4 pipeline demonstrates the sheer utility of the \textit{Alien Mathematics} paradigm. The AI is no longer a black box oracle; it is an explorer mapping out formal structures that human mathematicians can then rigorously audit, compile, and generalize.
